% =====================
% Chapter 12: Copying Data — Assignment, Shallow Copy, Deep Copy
% Audience: Architecture students new to coding
% Conventions: \h, \hh, \hhh, \code, NoteBox, minilab
% Code IDs: c12_*
% =====================

\h{Copying data: assignment, shallow copy, deep copy}
In an architectural office, making a copy is a daily task. But the type of copy you make is critical. Do you make a quick photocopy of a floor plan, or do you need to duplicate the entire project binder, including every referenced structural drawing and specification sheet? The first is a simple copy, but if a core structural drawing it refers to is updated, your simple copy is instantly out of date. The second is a complete, independent duplicate.

Python has a similar distinction when you want to duplicate data structures like lists. Sometimes, you don't get a true, independent copy; you get something that is still linked to the original. Understanding the difference between a shallow copy (the quick photocopy) and a deep copy (the full project binder duplication) is essential for avoiding bugs where you change one list and accidentally modify another.

Scripts often reuse “templates” of data: a room record, a list of levels, or a nested structure like a building -> floors -> rooms. Copying that data correctly avoids accidental edits that leak back into the original.

\begin{NoteBox}
\emph{Assignment} binds a new name to the \emph{same} object. 
A \emph{shallow copy} makes a new container but keeps references to the same inner objects. 
A \emph{deep copy} duplicates the container \emph{and} all nested objects.
\end{NoteBox}

\begin{tcolorbox}[% passing options here
    title=Focus of this chapter:,
    colframe=orange,
    colback = white,
%    colback=orange!40!yellow!30, % think of mixing 2 colors with sliders
    arc=2mm % sharp corners
 ]
\begin{itemize}
  \item Distinguish assignment vs. copying; use \c{is} and \c{id()} to check identity.
  \item Make shallow copies of lists, dicts, and sets safely.
  \item Decide when a deep copy is required for nested structures.
  \item Avoid common pitfalls: list replication, nested aliasing, default mutables.
\end{itemize}
\end{tcolorbox}

\hh{The Baseline: Assignment {=} is Not a Copy}
First, we must understand a fundamental concept. When you use the assignment operator (=) with a list, you are not creating a copy. You are simply creating a second name, or a second "label," that points to the exact same list in memory. Think of it as giving a coworker a key to your office. You both now have access to the one and only office. If your coworker moves a chair, the chair is moved for you too. That is, assignment does not copy—it creates another name pointing to the same object. Mutating
through either name changes the single shared object.

\code{c12_assign_is_alias}{python}{
    rooms = ["Lobby", "Cafe"]
    alias = rooms                  # assignment: no copy
    alias.append("Studio")         # mutate via alias
    print(rooms)                   # both names see the change
    print(rooms is alias)          # True: same object
    # Outcome:
    # ['Lobby', 'Cafe', 'Studio']
    # True
}{Assignment aliases the same object; no copy occurs.}

\code{c12_assign_alias_materials}{python}{
    # A list of primary structural materials
    original_materials = ["steel", "concrete", "wood"]

    # This does NOT create a new list. It just creates a new label (alias).
    shared_materials = original_materials

    # Modify the list using the 'shared_materials' label
    shared_materials.append("glass")

    print(f"Shared list: {shared_materials}")
    print(f"Original list: {original_materials}")  # The original has also changed!

    # Identity check: both names refer to the same object in memory
    print(shared_materials is original_materials)   # True
    # Outcome:
    # Shared list: ['steel', 'concrete', 'wood', 'glass']
    # Original list: ['steel', 'concrete', 'wood', 'glass']
    # True
}{Assignment binds a new name to the same list; mutations are shared.}

\hh{Shallow copy: new outer container, same inner objects}
A shallow copy creates a new list object, but it doesn't create copies of the items inside the list. Instead, it fills the new list with references to the original items. This is fine for simple lists of numbers or strings. The problem arises when you have a list that contains other lists (a nested list). A shallow copy will create a new outer list, but the inner lists it contains are still the exact same ones from the original.\par
Analogy: You photocopy a single-page project schedule. You have a new piece of paper, but the schedule lists tasks like "Phase A," "Phase B," etc. Those phases are complex projects themselves, detailed in other binders. Your photocopy still points to those same original binders. If someone changes something in the "Phase A" binder, the information your copy refers to has also changed.\par
Common ways: slicing, constructor call, or \c{.copy()}.

\code{c12_shallow_list_ways}{python}{
    base = [1, 2, 3]
    a = list(base)                 # constructor
    b = base[:]                    # full slice
    c = base.copy()                # method
    print(a == base, a is base)    # True False (same values, new object)
    print(b is base, c is base)    # False False
}{Three equivalent ways to shallow-copy a list.}

\code{c12_shallow_dict_set}{python}{
    info = {"name": "Lobby", "area": 24.0}
    cp1  = dict(info)              # constructor
    cp2  = info.copy()             # method
    print(cp1 == info, cp1 is info)  # True False
    levels = {1, 2, 3}
    cp3 = set(levels)
    cp4 = levels.copy()
    print(cp3 == levels, cp3 is levels)  # True False
}{Shallow copy for dicts and sets.}

\hhh{Shallow copy caveat: nested structures still share inner objects}

If the container holds other containers, only the outer shell is new. Inner pieces are
still shared.

\code{c12_shallow_nested_caveat}{python}{
    schedule = {
        "A": ["Lobby", "Cafe"],
        "B": ["Studio"]
    }
    shallow = schedule.copy()      # new dict, same inner lists
    shallow["A"].append("Gallery") # mutate inner list
    print(schedule["A"])           # original changed too
    print(schedule is shallow, schedule["A"] is shallow["A"])
    # Outcome:
    # ['Lobby', 'Cafe', 'Gallery']
    # False True
}{Shallow copy: new outer dict, shared inner lists.}

\code{c12_shallow_plan_alias}{python}{
    # A project plan, where each inner list is a phase
    original_plan = [
        ["Schematic Design", "Complete"],
        ["Design Development", "In Progress"]
    ]

    # Create a shallow copy (new outer list, same inner lists)
    shallow_copy_plan = original_plan.copy()

    # Modify a nested item IN THE COPY
    # Change the status of "Design Development"
    shallow_copy_plan[1][1] = "Complete"

    print(f"Shallow Copy: {shallow_copy_plan}")
    print(f"Original Plan: {original_plan}")  # The original has been accidentally modified!
    # Outcome:
    # Shallow Copy: [['Schematic Design', 'Complete'], ['Design Development', 'Complete']]
    # Original Plan: [['Schematic Design', 'Complete'], ['Design Development', 'Complete']]
}{Shallow copy: outer list is new, but inner lists are shared (aliased). Because the inner list ["Design Development", "In Progress"] was shared between both the original and the copy, changing it in one place changed it for both.}

\hh{Deep copy: duplicate everything recursively}
Deep copy is like duplicating the Entire Project Binder. A deep copy solves this problem. It creates a new list and then, recursively, creates brand new copies of all the items inside, no matter how deeply they are nested. This ensures that you have a completely independent duplicate of the original data structure.\par
Analogy: This is like duplicating the entire project binder. You don't just photocopy the top page; you go and find every single referenced drawing and specification sheet and make a fresh copy of those as well. Your new binder is a complete, standalone set of documents.\par
That is, deep copy is for nested structures you plan to mutate independently. Use the \c{copy} module.
\begin{NoteBox}
\textbf{Performance tip.} Deep copying large structures can be slow. Prefer shallow copies plus creating fresh inner
objects \emph{only} where you plan to mutate.
\end{NoteBox}

\code{c12_deepcopy_nested}{python}{
    import copy

    schedule = {
        "A": ["Lobby", "Cafe"],
        "B": ["Studio"]
    }
    deep = copy.deepcopy(schedule)  # full recursive copy
    deep["A"].append("Gallery")     # mutate copy
    print(schedule["A"])            # original unaffected
    print(schedule is deep, schedule["A"] is deep["A"])
    # Outcome:
    # ['Lobby', 'Cafe']
    # False False
}{Deep copy duplicates the container and all nested objects.}

\hh{When to use which}

\begin{itemize}
  \item \textbf{Assignment} when you truly want two names for the same object (e.g., pass-by-reference style).
  \item \textbf{Shallow copy} for flat containers or when you only need a new outer shell.
  \item \textbf{Deep copy} for nested containers you will modify independently.
\end{itemize}

\begin{NoteBox}
\textbf{Performance tip.} Deep copies can be expensive on large graphs. Prefer shallow copies plus
targeted creation of fresh inner objects where you plan to mutate.
\end{NoteBox}

\hh{Common copy patterns (and safe alternatives)}

\hhh{Grid pitfall: list replication aliases rows}

Using \c{*} to replicate a list of lists creates many references to \emph{the same} inner list.

\code{c12_grid_replication_pitfall}{python}{
    rows = 3
    cols = 3
    grid_bad = [[0] * cols] * rows   # replicated references!
    grid_bad[0][0] = 1
    print(grid_bad)                   # first column changed in all rows
    # Outcome:
    # [[1, 0, 0], [1, 0, 0], [1, 0, 0]]
}{List replication duplicates references, not rows.}

\code{c12_grid_comprehension_ok}{python}{
    rows = 3
    cols = 3
    grid_ok = [[0 for _ in range(cols)] for _ in range(rows)]  # fresh row each time
    grid_ok[0][0] = 1
    print(grid_ok)                         # only one cell changes
    # Outcome:
    # [[1, 0, 0], [0, 0, 0], [0, 0, 0]]
}{Use a comprehension to create independent inner lists.}

\hhh{Copying records: dict-of-lists}

Decide shallow vs deep based on whether inner lists will be mutated.

\code{c12_copy_records_compare}{python}{
    import copy

    template = {
        "floor": 1,
        "tags": ["public", "egress"],
    }

    # Shallow copy: safe if you won't mutate 'tags'
    a = template.copy()
    a["floor"] = 2                 # independent (new int)
    a["tags"].append("display")    # mutates shared list!

    # Deep copy: safe if you will mutate nested lists
    b = copy.deepcopy(template)
    b["tags"].append("gallery")    # independent

    print(template["tags"])        # shared 'display' visible here
    print(b["tags"])               # independent 'gallery'
    # Outcome:
    # ['public', 'egress', 'display']
    # ['public', 'egress', 'gallery']
}{Choose deep copy if nested lists/dicts will change.}

\hh{The {copy} module and custom objects}

\begin{itemize}
  \item \c{copy.copy(obj)} makes a shallow copy; \c{copy.deepcopy(obj)} makes a deep copy.
  \item For your own classes, you can customize copying with \c{\_\_copy\_\_} and \c{\_\_deepcopy\_\_} methods (advanced).
\end{itemize}

\code{c12_copy_module_demo}{python}{
    import copy

    class Room:
        def __init__(self, name, tags):
            self.name = name
            self.tags = list(tags)

    r1 = Room("Lobby", ["public"])
    r2 = copy.copy(r1)             # shallow: new Room, same inner 'tags' list ref
    r3 = copy.deepcopy(r1)         # deep: new Room, new 'tags' list

    r2.tags.append("display")
    r3.tags.append("gallery")

    print(r1.tags)                 # shared with r2
    print(r3.tags)                 # independent
    # Outcome:
    # ['public', 'display']
    # ['public', 'gallery']
}{Shallow vs deep copy for custom objects with inner mutables.}

\hh{Quick reference}

\renewcommand{\arraystretch}{1.15}
\begin{table}[h]
\centering
\begin{tabularx}{\textwidth}{|l|X|X|}
\hline
\textbf{Type} & \textbf{Shallow copy} & \textbf{Notes} \\ \hline
list & \texttt{list(a)}, \texttt{a[:]}, \texttt{a.copy()} & Nested lists remain shared; use \texttt{deepcopy} or build fresh inner lists. \\ \hline
dict & \texttt{dict(d)}, \texttt{d.copy()} & Nested values remain shared. \\ \hline
set & \texttt{set(s)}, \texttt{s.copy()} & Elements are shared references. \\ \hline
tuple & same object (immutable) & Usually no copy needed; contains references. \\ \hline
custom & \texttt{copy.copy(obj)} & Deep copy with \texttt{copy.deepcopy(obj)} if nested mutables must be independent. \\ \hline
\end{tabularx}
\caption{Shallow copy recipes and cautions.}
\end{table}

\hh{Copy behavior by data type}

Not all data needs the same kind of copy. Use this guide to decide quickly (see Table \ref{tab:copying_Data_Reference}):
\begin{itemize}
  \item \textbf{Immutable types} (numbers, strings, tuples of immutables) are safe: assignment and copies behave the same because you cannot mutate the value.
  \item \textbf{Mutable containers} (lists, dicts, sets) need care: shallow copies share inner objects; deep copies duplicate nested content.
  \item \textbf{Mixed/nested} data (e.g., tuple containing a list) behaves like its mutable parts.
\end{itemize}

\hhh{Immutable scalars and strings}

Assignment and copy give you effectively the same result; there is nothing to mutate.

\code{c12_copy_immutable}{python}{
    import copy

    a = 42
    b = a                  # assignment
    c = copy.copy(a)       # shallow copy
    d = copy.deepcopy(a)   # deep copy
    print(a == b, a is b)  # True True (may be same object)
    print(a is c, a is d)  # True True

    s1 = "Lobby"
    s2 = s1
    s3 = copy.copy(s1)
    s4 = copy.deepcopy(s1)
    print(s1 == s3, s1 is s3)  # True True
}{Immutables: assignment/copy/deepcopy all safe; mutation is impossible.}

\hhh{Tuples: immutable shell, but contents may be mutable}

A tuple itself is immutable; a \emph{shallow copy} of a tuple is the same object. If it holds a mutable item, that inner object is shared unless you deep copy.

\code{c12_copy_tuple_with_list}{python}{
    import copy

    t = (1, [2, 3])
    t_shallow = copy.copy(t)      # same tuple object
    t_deep = copy.deepcopy(t)     # new tuple; inner list duplicated

    t[1].append(4)                # mutate the inner list
    print(t)                      # (1, [2, 3, 4])
    print(t_shallow)              # shares inner list -> (1, [2, 3, 4])
    print(t_deep)                 # independent -> (1, [2, 3])
    print(t is t_shallow, t is t_deep)  # True False
}{Tuples are immutable, but inner mutables require deep copy for independence.}

\hhh{Lists: classic mutable container}

Shallow copies make a new outer list but share inner objects; deep copies duplicate everything.

\code{c12_copy_list_nested}{python}{
    import copy

    plan = [["SD", "Complete"], ["DD", "In Progress"]]
    plan_shallow = plan[:]             # shallow
    plan_deep = copy.deepcopy(plan)    # deep

    plan_shallow[1][1] = "Complete"    # edit nested item in shallow copy
    print(plan)                        # nested change leaks back
    print(plan_deep)                   # deep copy stays independent
    # Outcome:
    # [['SD', 'Complete'], ['DD', 'Complete']]
    # [['SD', 'Complete'], ['DD', 'In Progress']]
}{Lists: shallow shares nested items; deep keeps them separate.}

\hhh{Dicts and sets}

Same rule as lists: shallow copies share nested values; deep copies duplicate them.

\code{c12_copy_dict_set}{python}{
    import copy

    info = {"name": "Studio", "tags": ["public"]}
    info_shallow = info.copy()
    info_deep = copy.deepcopy(info)

    info_shallow["tags"].append("display")
    print(info["tags"])       # ['public', 'display']  (shared)
    print(info_deep["tags"])  # ['public']             (independent)

    floors = {1, 2, 3}
    floors_copy = floors.copy()     # shallow copy is enough (no nesting)
    print(floors is floors_copy)    # False (new set object)
}{Dicts/sets: shallow shares nested content; deep copies it. Flat sets need only shallow.}

\hhh{Bytes vs bytearray (immutable vs mutable)}

\code{c12_copy_bytes_bytearray}{python}{
    import copy

    b = b"ABC"                     # bytes: immutable
    ba = bytearray(b"ABC")         # bytearray: mutable

    print(copy.copy(b) is b)       # True
    print(copy.deepcopy(b) is b)   # True

    ba_copy = copy.copy(ba)
    ba_copy[0] = ord(b"Z")
    print(ba, ba_copy)             # bytearrays are independent objects
}{Immutables (bytes) return the same object; mutables (bytearray) copy as containers.}

\hhh{Custom objects with inner mutables}

`copy.copy(obj)` makes a new instance but keeps references to inner mutables; `copy.deepcopy(obj)` duplicates those too.

\code{c12_copy_custom}{python}{
    import copy

    class Room:
        def __init__(self, name, tags):
            self.name = name
            self.tags = list(tags)   # mutable attribute

    r1 = Room("Lobby", ["public"])
    r2 = copy.copy(r1)               # shallow: shares 'tags' list
    r3 = copy.deepcopy(r1)           # deep: new 'tags' list

    r2.tags.append("display")
    r3.tags.append("gallery")
    print(r1.tags)                   # ['public', 'display']  (shared with r2)
    print(r3.tags)                   # ['public', 'gallery']  (independent)
}{For classes with mutable attributes, use deepcopy for independent nested state.}

\hh{Quick reference by type}

\renewcommand{\arraystretch}{1.12}
\begin{table}[h]
\centering
\begin{tabularx}{\textwidth}{|l|X|X|}
\hline
\textbf{Type} & \textbf{Assignment} & \textbf{Shallow vs deep copy} \\ \hline
int, float, bool, str & New name points to same immutable value; safe & \texttt{copy} / \texttt{deepcopy} usually return the same object; no mutation possible \\ \hline
tuple (all immutables) & Alias is fine & Shallow = same tuple; deep = same effect; nothing to copy \\ \hline
tuple (has mutables) & Alias shares inner mutables & Shallow shares inner objects; \textbf{deep} duplicates inner mutables \\ \hline
list & Alias shares list; mutations visible via all names & Shallow: new list, shared items; Deep: new list and new nested items \\ \hline
dict & Alias shares mapping; mutations visible via all names & Shallow: new dict, shared values; Deep: duplicates nested values \\ \hline
set & Alias shares set; mutations shared & Shallow: new set (elements referenced); Deep: matters only if elements are mutable containers \\ \hline
bytes & Immutable; alias safe & Copy/deepcopy return same object \\ \hline
bytearray & Mutable; alias shares buffer & Shallow makes new bytearray; Deep same as shallow \\ \hline
custom class (with mutables) & Alias shares entire object & Shallow: new instance, shared inner mutables; Deep: duplicates inner mutables \\ \hline
\end{tabularx}
\caption{How assignment, shallow copy, and deep copy behave for common types.\label{tab:copying_Data_Reference}}
\end{table}

\begin{NoteBox}
\textbf{Library note.} Third-party containers (e.g., NumPy arrays, pandas DataFrames) have their \emph{own} copying rules (views vs copies). Always check their docs; built-in Python rules above may not apply.
\end{NoteBox}

\hh{Mini-Lab A — Safe room grid}

Create a 4x3 grid of zeros. Set cell (row 0, col 0) to 1. Show that only that cell changes.

\code{c12_lab_grid}{python}{
    rows, cols = 4, 3
    grid = [[0 for _ in range(cols)] for _ in range(rows)]
    grid[0][0] = 1
    print(grid)
    # Outcome:
    # [[1, 0, 0], [0, 0, 0], [0, 0, 0], [0, 0, 0]]
}{Build independent rows using a comprehension.}

\hh{Mini-Lab B — Duplicate a project template}

Given a nested template, make an editable duplicate that will not affect the original.

\code{c12_lab_project_template}{python}{
    import copy

    template = {
        "meta": {"name": "Studio A", "level": 2},
        "rooms": [
            {"name": "Lobby", "tags": ["public"]},
            {"name": "Cafe",  "tags": ["service"]},
        ],
    }

    dup = copy.deepcopy(template)      # safe duplication
    dup["meta"]["name"] = "Studio B"
    dup["rooms"][0]["tags"].append("display")

    print(template["meta"]["name"])    # unchanged
    print(template["rooms"][0]["tags"])# unchanged
    print(dup["meta"]["name"])         # changed
    print(dup["rooms"][0]["tags"])     # changed only in copy
    # Outcome:
    # Studio A
    # ['public']
    # Studio B
    # ['public', 'display']
}{Use deepcopy to edit a nested structure safely.}

\hh{Mini-Lab C — Shallow copy where it’s enough}

Copy a flat list of floor numbers, edit the copy, and show the original is unchanged.

\code{c12_lab_flat_copy}{python}{
    floors = [1, 2, 3, 4]
    copy_f = floors[:]                # shallow copy is fine for a flat list
    copy_f.append(5)
    print(floors)                     # original unchanged
    print(copy_f)                     # copy has 5
    # Outcome:
    # [1, 2, 3, 4]
    # [1, 2, 3, 4, 5]
}{Shallow copy is sufficient for flat, independent edits.}

\hh{Key Takeaways: Assignment vs. Copying}
\begin{tcolorbox}[
    colframe=orange,
    colback=white,
    arc=2mm
]
\begin{itemize}
    \item \textbf{Assignment ({b = a})}: This is not a copy. It's just a second name or label pointing to the exact same object in memory. Modifying \c{b} will also modify \c{a}.
    \item \textbf{Shallow Copy ({b = copy.copy(a)})}: Creates a new top-level object, but fills it with \emph{references} to the items in the original. It is safe for simple, "flat" lists (like a list of numbers), but changes to nested objects will affect both copies.
    \item \textbf{Deep copy ({b = copy.deepcopy(a)})}: Creates a completely new, fully independent duplicate of the original object and everything inside it, no matter how deeply nested. This is the safest way to ensure two variables do not interfere with each other.
    \item \textbf{Import first:} To use shallow or deep copy, you must first import the library with \c{import copy}.
    \item \textbf{When to use:} Use a deep copy whenever you need to modify a complex data structure (like a list of lists) without affecting the original source data.
\end{itemize}
\end{tcolorbox}

\hh{Self-check}
\begin{itemize}
  \item Explain why \texttt{a = b} is not a copy. How do you test whether two names share an object?
  \item Show a case where \texttt{dict.copy()} is not enough and \texttt{deepcopy} fixes it.
  \item Rewrite a grid built with \texttt{[[0]*n]*m} to avoid aliasing.
  \item When is deep copying overkill? Give one alternative.
\end{itemize}
